\section{Phase 1: Separability-Guided Feature Configuration}
\label{sec:phase1}
\label{sec:phase1_method}

\subsection{Motivation}
\label{sec:phase1_motivation}

The performance of handcrafted feature-based classification systems depends
heavily on the parameter configuration of the underlying feature extractors.
Parameters such as quantization levels in the GLCM, neighborhood settings in
LBP, decomposition levels in the DWT, and cell structures in HOG directly
influence the discriminative quality of the extracted representations.

Most existing studies either employ default parameter values or optimize
descriptors using classifier-dependent criteria such as classification
accuracy. Although effective, classifier-dependent optimization introduces a
coupling between feature extraction and classifier behavior, making it difficult
to determine whether performance improvements arise from the quality of the
feature representation or from the characteristics of a particular learning
algorithm. To overcome this limitation, this work adopts a
classifier-independent optimization strategy based on feature separability:
candidate configurations are assessed according to their ability to separate
tumor classes in the feature space before any classifier training is performed.
The objective of Phase~1 is therefore to identify descriptor configurations that
maximize class discrimination while remaining independent of any specific
machine learning model.

\subsection{Descriptor Search Space}
\label{sec:phase1_search}

Four handcrafted feature families were investigated. For each, a grid of
candidate parameter configurations was generated from the training images and
evaluated, yielding $144$ GLCM, $36$ LBP, $21$ DWT, and $12$ HOG configurations
($213$ in total).

\subsubsection{GLCM Search Space}
The Gray-Level Co-occurrence Matrix describes second-order texture statistics by
modeling spatial relationships between pixel intensities. The explored
parameters were the number of gray levels, the set of pixel distances, the set
of angular directions, and matrix symmetrization. For each configuration, six
texture properties---contrast, dissimilarity, homogeneity, energy, correlation,
and angular second moment (ASM)---were extracted and averaged over all
distance--angle pairs, producing a six-dimensional descriptor.

\subsubsection{LBP Search Space}
Local Binary Patterns characterize local texture structures by comparing
neighboring pixels with a central pixel. The explored parameters were the number
of sampling points $P$, the neighborhood radius $R$, and the encoding method
(\texttt{uniform}, \texttt{nri\_uniform}, and \texttt{ror}). The resulting
pattern histograms were used as texture descriptors for separability evaluation.

\subsubsection{DWT Search Space}
The Discrete Wavelet Transform provides a multi-resolution representation of
image content by decomposing an image into frequency sub-bands. The search space
comprised seven wavelet families and three decomposition levels. Four
statistical measures were extracted from each sub-band (approximation and
detail coefficients), yielding $4(3L+1)$ features for a decomposition of level
$L$.

\subsubsection{HOG Search Space}
The Histogram of Oriented Gradients captures local shape and structural
information through gradient orientation distributions. The explored parameters
were the number of orientation bins, the cell size (pixels per cell), and the
block size (cells per block); block normalization was fixed to L2-Hys. Each
combination generated a different gradient-based representation
(Figure~\ref{fig:hog_schema}).

\begin{figure}[t]
\centering
\includegraphics[width=\linewidth]{hog_schema}
\caption{Illustration of HOG descriptor computation on a preprocessed MRI slice.}
\label{fig:hog_schema}
\end{figure}

\subsection{Separability Metrics}
\label{sec:separability}

To assess descriptor quality independently of classifier performance, three
complementary separability metrics were employed. Prior to evaluation, each
candidate feature matrix was standardized, and all metrics were computed on the
training set only.

\subsubsection{Fisher Discriminant Ratio (FDR)}
The Fisher Discriminant Ratio relates inter-class separation to intra-class
variability. For a single feature it is defined as
\begin{equation}
\mathrm{FDR} =
\frac{\sum_{c=1}^{C} n_c (\mu_c - \mu)^2}
     {\sum_{c=1}^{C} n_c \sigma_c^2},
\end{equation}
where $C$ is the number of classes, $n_c$, $\mu_c$, and $\sigma_c^2$ are the
size, mean, and variance of class $c$, and $\mu$ is the global mean. The score
is averaged across all features; higher values indicate better class
discrimination.

\subsubsection{Mutual Information (MI)}
Mutual Information measures the statistical dependency between a feature $X$ and
the class label $Y$,
\begin{equation}
\mathrm{MI}(X,Y) = \sum_{x \in X}\sum_{y \in Y}
p(x,y)\,\log \frac{p(x,y)}{p(x)\,p(y)}.
\end{equation}
In practice it is estimated with a nearest-neighbor estimator
($k=5$) and averaged across features; higher values indicate more
class-relevant information.

\subsubsection{Davies--Bouldin Index (DBI)}
The Davies--Bouldin Index~\cite{davies1979} evaluates cluster compactness and
separation,
\begin{equation}
\mathrm{DBI} = \frac{1}{C}\sum_{i=1}^{C}
\max_{j \neq i}\left(\frac{S_i + S_j}{M_{ij}}\right),
\end{equation}
where $S_i$ and $S_j$ denote intra-cluster dispersion and $M_{ij}$ the distance
between the centroids of classes $i$ and $j$. To keep the computation efficient
in high dimensions, the DBI is evaluated on a PCA-reduced projection (at most
$20$ components) of the standardized features. Unlike FDR and MI, lower DBI
values indicate better separability.

\subsection{Composite Separability Score}
\label{sec:phase1_composite}

Because each metric captures a different aspect of class discrimination, a
composite score was constructed as a unified ranking criterion. Within each
descriptor family, metric values were first normalized to $[0,1]$ across the
candidate configurations. For FDR and MI,
\begin{equation}
M_{\mathrm{norm}} = \frac{M - M_{\min}}{M_{\max} - M_{\min}},
\end{equation}
and for the DBI, where lower values are preferred,
\begin{equation}
\mathrm{DBI}_{\mathrm{norm}} = 1 - \frac{\mathrm{DBI} - \mathrm{DBI}_{\min}}
{\mathrm{DBI}_{\max} - \mathrm{DBI}_{\min}}.
\end{equation}
The composite score is the equally weighted average
\begin{equation}
\begin{split}
\mathrm{Score} = w_1\,\mathrm{FDR}_{\mathrm{norm}}
+ w_2\,\mathrm{MI}_{\mathrm{norm}}
+ w_3\,\mathrm{DBI}_{\mathrm{norm}}, \\
w_1 = w_2 = w_3 = \tfrac{1}{3}.
\end{split}
\end{equation}
Because normalization is performed \emph{within} each descriptor family, the
composite score ranks configurations of the same descriptor against one another;
it should not be interpreted as an absolute, cross-descriptor measure of feature
quality.

\subsection{Optimization Procedure}
\label{sec:phase1_procedure}

The optimization procedure was performed independently for each descriptor
family. For every candidate configuration, features were extracted from the
training images; the FDR, MI, and DBI were computed; the metric values were
normalized and combined into the composite score; and the configurations were
ranked. The highest-ranked configuration for each family was retained for the
subsequent experiments. The procedure can be summarized as: configuration
generation $\rightarrow$ feature extraction $\rightarrow$ separability
evaluation $\rightarrow$ composite scoring $\rightarrow$ ranking $\rightarrow$
optimal configuration selection.

Figure~\ref{fig:phase1_search} visualizes the composite separability score across
the explored parameter grids for all four descriptor families.

\begin{figure*}[t]
\centering
\begin{subfigure}[b]{0.48\textwidth}
\includegraphics[width=\linewidth]{glcm_param_search}
\caption{GLCM}
\end{subfigure}
\hfill
\begin{subfigure}[b]{0.48\textwidth}
\includegraphics[width=\linewidth]{lbp_param_search}
\caption{LBP}
\end{subfigure}\\[0.5em]
\begin{subfigure}[b]{0.48\textwidth}
\includegraphics[width=\linewidth]{dwt_param_search}
\caption{DWT}
\end{subfigure}
\hfill
\begin{subfigure}[b]{0.48\textwidth}
\includegraphics[width=\linewidth]{hog_param_search}
\caption{HOG}
\end{subfigure}
\caption{Phase~1 parameter search: composite separability score across candidate
configurations for each descriptor family.}
\label{fig:phase1_search}
\end{figure*}

\subsection{Optimal Descriptor Configurations}
\label{sec:phase1_optimal}

Table~\ref{tab:phase1} reports the selected configuration for each descriptor
family, together with its dimensionality and composite score. The best GLCM
configuration combined four pixel distances at a single orientation with $128$
gray levels and an asymmetric co-occurrence matrix. The optimal LBP descriptor
used twelve uniform sampling points at radius one. For the DWT, the
\texttt{db2} wavelet at a single decomposition level offered the best
separability score. The optimal HOG configuration used six orientation bins with
$16\times16$-pixel cells and $3\times3$-cell blocks.

\begin{table}[t]
\centering
\caption{Optimal descriptor configurations selected in Phase~1. Composite scores
are normalized within each descriptor family and are therefore not directly
comparable across families.}
\label{tab:phase1}
\renewcommand{\arraystretch}{1.3}
\begin{tabular}{@{}llcc@{}}
\toprule
\textbf{Descriptor} & \textbf{Optimal parameters} & \textbf{Dim.} & \textbf{Score} \\
\midrule
GLCM & dist.\ $\{1,2,4,8\}$, $1$ angle ($0^{\circ}$), & $6$ & $0.7468$ \\
     & $128$ levels, asymmetric & & \\
LBP  & $P=12$, $R=1$, uniform & $14$ & $0.9906$ \\
DWT  & \texttt{db2}, level $1$ & $16$ & $0.9837$ \\
HOG  & $6$ orient., $16\times16$ cell, & $1176$ & $0.9898$ \\
     & $3\times3$ block, L2-Hys & & \\
\bottomrule
\end{tabular}
\end{table}

Figure~\ref{fig:phase1_best} plots the separability metrics for the selected
optimal configuration within each family.

\begin{figure*}[t]
\centering
\begin{subfigure}[b]{0.48\textwidth}
\includegraphics[width=\linewidth]{glcm_best_params}
\caption{GLCM}
\end{subfigure}
\hfill
\begin{subfigure}[b]{0.48\textwidth}
\includegraphics[width=\linewidth]{lbp_best_params}
\caption{LBP}
\end{subfigure}\\[0.5em]
\begin{subfigure}[b]{0.48\textwidth}
\includegraphics[width=\linewidth]{dwt_best_params}
\caption{DWT}
\end{subfigure}
\hfill
\begin{subfigure}[b]{0.48\textwidth}
\includegraphics[width=\linewidth]{hog_param_search}
\caption{HOG (search grid)}
\end{subfigure}
\caption{Separability profiles for the optimal GLCM, LBP, and DWT configurations
and the HOG parameter grid.}
\label{fig:phase1_best}
\end{figure*}

\subsection{Discussion of Phase 1 Results}
\label{sec:phase1_discussion}

The optimization results demonstrate that descriptor parameterization has a
substantial impact on feature discriminability. Marked differences were observed
between candidate configurations within each descriptor family, confirming that
default parameter settings do not necessarily produce optimal representations
for brain tumor MRI classification.

The proposed separability-guided framework offers several advantages. First, it
enables an objective evaluation of descriptor quality without relying on
classifier performance. Second, it reduces the risk of overfitting feature
configurations to a specific learning algorithm. Third, it provides a systematic
methodology for selecting feature extractors before classifier benchmarking is
performed. It should be emphasized that the composite scores are a
classifier-independent \emph{selection} criterion computed on the training set,
rather than an estimate of classification accuracy; the predictive value of each
representation is assessed separately in Phase~2.

The optimized descriptors obtained in this phase constitute the foundation for
the classifier benchmarking experiments presented in the next section. By
ensuring that all classifiers operate on carefully selected feature
representations, the subsequent comparison focuses on differences attributable
to the classifiers themselves rather than to variations in feature quality.
